\begin{afterword}
\summaryhead

The post asks why humans, made by a process that favours only inclusive genetic fitness, do not want fitness itself. An agent that did is possible in principle, the author argues: a mind born with the concept of fitness and a library of explicit strategies would eat, explore and have sex only as far as these served reproduction, and it would not invent condoms. Two objections, about hunger and curiosity, are answered to show that such an agent needs no separate drives.

Evolution did not build one, the author says, because evolution cannot redesign a system all at once. For millions of years brains learned from rewards, such as the taste of fat or the pleasure of sex, that correlated with reproduction only at long range. When general reasoning appeared, it was put to work getting those existing rewards, which was far simpler than rebuilding the mind around fitness. So evolution's single goal splintered into ``a thousand shards of desire.'' The author judges this a good thing, since the alternative is a future of pure replicators, and answers the objection that our tastes also come from evolution with ``So what?''

\responsehead

The picture of people as pursuing the rewards evolution installed, not fitness itself, is the one the book has already given in ``Adaptation-Executers, not Fitness-Maximizers,'' with credit to Tooby and Cosmides. This post adds a question in a programmer's terms, why evolution did not build an explicit fitness maximizer, and answers it plainly: selection cannot rewrite a whole design, so reasoning was attached to the rewards that were already there. The distinction between rewards tied to fitness at long range and to behaviour at short range is the useful part, and the post's example, an apple that simply tastes good, makes it concrete.

The central step is an inference presented as history. That early reasoners ``reasoned consequentially about how to get the existing reinforcers'' because this was ``a relatively simple hack'' is a claim about what happened in hominid evolution. The post supports it only by what selection could easily find. The inference fits the earlier posts and may well be right, but the post gives no evidence of the kind a claim about the past would need.

The ending is a value judgment, and the post marks it as one (``I think,'' ``though I'm a human who says so''). The alternative it weighs is descendants who pursue fitness alone. It then states the obvious objection, that the tastes doing the judging were also made by evolution, and answers ``So what?'' That is a position, that the origin of a value does not settle whether to keep it. The post gives no argument for it here.

In short: a plain answer to why human desires are not a desire for fitness, whose key historical step is an inference stated as fact, and whose closing judgment answers its own objection with ``So what?''
\end{afterword}
